% Reading companion to quirks_p1p2_note.pdf (v0.4): what Kang-Luty, Knapen et al. and Li et al. wrote,
% what P1 changes, and where P1 enters their work.
% Build (from theory_track/paper/companion/):  xelatex prior_work_companion && xelatex prior_work_companion
\documentclass[11pt]{article}
% Preamble for the C001 quirk P1/P2 note (xelatex + bibtex/natbib).
% Adapted from jfc/src/conventions/preamble.tex (graphics/float/caption/breaking settings kept as is);
% additions: fonts for xelatex, maths, natbib, hyperref, notation macros, uncertainty tags, TODO marker.

% --- Fonts / encoding (xelatex) ---
\usepackage{fontspec}
\usepackage{amsmath,amssymb}
\usepackage[a4paper,hmargin=2.0cm,vmargin=1.8cm]{geometry}   % vmargin 2.0 -> 1.8 cm: page budget (fix R2)
\usepackage{microtype}

% --- Graphics and floats (JFC) ---
\usepackage{graphicx}
\usepackage{placeins}
\usepackage{float}
\usepackage{needspace}
\setkeys{Gin}{height=0.45\linewidth,keepaspectratio}

% --- Caption formatting (JFC) ---
\usepackage[font=small,labelfont=bf]{caption}

% --- Tables ---
\usepackage{booktabs}
\usepackage{tabularx}
\usepackage{longtable}
\setlength{\LTpre}{0pt}
\setlength{\LTpost}{0pt}

% --- Float placement (JFC, relaxed) ---
\renewcommand{\topfraction}{0.9}
\renewcommand{\bottomfraction}{0.9}
\renewcommand{\textfraction}{0.05}
\renewcommand{\floatpagefraction}{0.6}
\setcounter{topnumber}{3}
\setcounter{bottomnumber}{3}
\setcounter{totalnumber}{6}

% --- Paragraph and page breaking (JFC) ---
\widowpenalty=10000
\clubpenalty=10000
\raggedbottom

% --- References, links ---
\usepackage[numbers,sort&compress]{natbib}
\usepackage{xcolor}
\usepackage[colorlinks=true,linkcolor=blue!50!black,citecolor=green!35!black,urlcolor=blue!50!black]{hyperref}
\usepackage{enumitem}
\usepackage{multicol}   % two-column bibliography (page budget, fix R2)
\setlist{itemsep=1pt,topsep=3pt}

% --- Units ---
\newcommand{\eV}{\,\mathrm{eV}}
\newcommand{\keV}{\,\mathrm{keV}}
\newcommand{\MeV}{\,\mathrm{MeV}}
\newcommand{\GeV}{\,\mathrm{GeV}}
\newcommand{\TeV}{\,\mathrm{TeV}}
\newcommand{\um}{\,\mu\mathrm{m}}
\newcommand{\mm}{\,\mathrm{mm}}
\newcommand{\cm}{\,\mathrm{cm}}
\newcommand{\ns}{\,\mathrm{ns}}
\newcommand{\ps}{\,\mathrm{ps}}
\newcommand{\fb}{\,\mathrm{fb}}
\newcommand{\abinv}{\,\mathrm{ab}^{-1}}
\newcommand{\tesla}{\,\mathrm{T}}

% --- Notation (one symbol per quantity; table in Appendix A) ---
\newcommand{\mQ}{m_Q}
\newcommand{\Lam}{\Lambda}
\newcommand{\Ds}{\Delta s}
\newcommand{\DsB}{\Delta s_B}
\newcommand{\DsMS}{\Delta s_{\rm MS}}
\newcommand{\DsBM}{\Delta s_{BM}}
\newcommand{\epsB}{\varepsilon_B}
\newcommand{\epsp}{\varepsilon_\perp}
\newcommand{\kB}{k_B}
\newcommand{\gst}{\gamma^*}
\newcommand{\bst}{\beta^*}
\newcommand{\lmax}{\ell_{\max}}
\newcommand{\lzero}{\ell_0}
\newcommand{\bodd}{b_{\rm odd}}
\newcommand{\rot}{\vartheta}
\newcommand{\Nosc}{N_{\rm osc}}
\newcommand{\RIT}{R_{\rm IT}}
\newcommand{\Bzero}{B_0}
\newcommand{\Estar}{\mathbf{E}^*}
\newcommand{\Bstar}{\mathbf{B}^*}
\newcommand{\Lcross}{\Lambda_{\rm cross}}
\newcommand{\sepIT}{d_{\rm IT}}
\newcommand{\Eff}{\mathcal{E}}
\newcommand{\Nbkg}{N_{\rm bkg}}
\newcommand{\Nbkgeightyfour}{N_{\rm bkg}^{84}}
\newcommand{\Ntrig}{N_{\rm trig}}
\newcommand{\Tseed}{T_{\rm seed}}
\newcommand{\wcore}{w_{\rm core}}
\newcommand{\chicore}{\chi^2_{\rm core}}
\newcommand{\dninetyfive}{\delta_{95}}
\newcommand{\Etil}{\tilde{E}}

% --- Uncertainty tags (as in the stage reports) ---
\newcommand{\EST}{\textsc{estimate}}
\newcommand{\UNV}{\textsc{unverified}}
\newcommand{\VER}{\textsc{verified}}

% --- Open-item marker for Phase B (must be empty at the final compile) ---
\newcommand{\TODO}[1]{{\color{red!75!black}\textsf{[TODO (Phase B): #1]}}}

% --- Stage-report pointer (appendices point to the stage reports for full detail) ---
\usepackage{xurl}
\newcommand{\rep}[1]{{\small\nolinkurl{#1}}}
\newcommand{\cmd}[1]{{\small\ttfamily #1}}
\emergencystretch=2em

% --- Compact table of contents (page budget, fix R2): less space between section entries ---
\usepackage{etoolbox}
\makeatletter
\patchcmd{\l@section}{\addvspace{1.0em \@plus\p@}}{\addvspace{0.25em \@plus\p@}}{}{\typeout{TOC patch failed}}
\makeatother

\usepackage{mdframed}
\usepackage{array}

% quoted source text: grey bar on the left
\newmdenv[topline=false,bottomline=false,rightline=false,linewidth=2.2pt,linecolor=black!30,
  leftmargin=0pt,rightmargin=0pt,innerleftmargin=9pt,innerrightmargin=0pt,innertopmargin=4pt,
  innerbottommargin=4pt,skipabove=6pt,skipbelow=6pt,backgroundcolor=black!3]{srcquote}
% our assessment: coloured bar
\newmdenv[topline=false,bottomline=false,rightline=false,linewidth=2.2pt,linecolor=blue!50!black,
  leftmargin=0pt,rightmargin=0pt,innerleftmargin=9pt,innerrightmargin=0pt,innertopmargin=4pt,
  innerbottommargin=4pt,skipabove=6pt,skipbelow=6pt]{verdict}
\newcommand{\src}[1]{\par\smallskip\noindent{\footnotesize\textsf{#1}}}
\newcommand{\vhead}[1]{\noindent\textbf{\textsf{#1}}\par\smallskip}
\newcommand{\note}[1]{{\small(note #1)}}

\title{Prior work on the quirk plane thickness, and where P1 enters it\\[3pt]
\large A reading companion to Kang--Luty, Knapen et al.\ and Li et al.}
\author{Companion to \emph{quirks\_p1p2\_note.pdf}, v0.4}
\date{19 October 2026}

\begin{document}
\maketitle

\begin{center}
\fbox{\parbox{0.92\linewidth}{\small\textbf{How to read this.} Grey boxes are text quoted from the papers. The quotes were
taken from the arXiv PDFs cached in \rep{research/cache/pdf/} (0805.4642, 1708.02243, 1911.02223, 2002.07503, 1811.08903).
The equations in grey boxes were transcribed from the PDF text layer and retyped; compare them with your printed copy.
Equation tags such as (KL~5.64) or (Li20~A.17) give the source paper's own numbers. ``Note~(7)'' means Eq.~(7) of our
note. Blue boxes are our assessment. Statements tagged \EST{} are our own reading or arithmetic in this companion; they are not
in the note and have not been reviewed.}}
\end{center}

\setcounter{tocdepth}{1}
\tableofcontents

%======================================================================================================
\section{Short answer}\label{sec:short}

None of the four papers makes a large error, and P1 changes none of their published limits: every one of them obtained the
plane thickness used in its collider analysis from a numerical simulation, not from the analytic estimates discussed here.
What P1 adds is the analytic origin of a quantity that all four touched but none derived: how far the solenoid field
tilts the quirk plane. Measured against P1, each paper has one or two statements that are approximate, and one
statement is misleading:

\begin{itemize}
\item \textbf{Kang \& Luty} estimate the field-induced separation ``after one oscillation'' as $\Delta r\sim(v_{\rm cm}B/m_Q)(m_Q/\Lambda^2)^2$.
  This is correct in scale. It is the miss distance $\bodd$ of P1 (note~(5)), and it exists only at odd crossings.
\item \textbf{Knapen et al.} estimate the plane rotation per oscillation as $\Delta\phi\propto B/\Lambda^2$, independent of $m_Q$.
  The scale $q|\Estar|/\sigma$ and the $m_Q$ independence are right. The estimate describes the back-and-forth swing within one period and
  misses the coherent, logarithmically enhanced rotation at each odd crossing (note~(7)). The true scaling is
  $B\ln(1/B)$, giving $\DsB\approx B^{0.9}\Lambda^{-1.8}$, not $B\Lambda^{-2}$.
\item \textbf{Li et al.\ 2019}, Eq.~(III.12), give the plane thickness as $d_{\min}\propto m_Q B/\Lambda^4$.
  \emph{This is the misleading one.} That formula has the scaling of the miss distance, not of the accumulated plane-fit
  thickness. It falls below the real thickness by $\times6$ at 2\,keV, $\times15$ at 3\,keV and $\times200$ at 10\,keV, and it predicts a mass dependence that their own Fig.~4 does not show.
\item \textbf{Li et al.\ 2020} have the right structure: a small rotation $\eta$ of the system per period, which adds up to the
  thickness. However, they determine $\eta$ numerically and assume it linear in the field, $r|\eta|\sim d_E$ (Eq.~A.17). P1 gives the
  analytic $\eta$ (note~(7)); it carries $\ln(\lmax/\bodd)\approx8$--14, so it is not linear in $B$. Their ``linear in the charge'' holds to
  about 7\% for $|q|=2$.
\end{itemize}

Table~\ref{tab:summary} at the end collects the verdicts.

%======================================================================================================
\section{The P1 chain in four steps}\label{sec:chain}

The note's Section~3 is a chain of four steps. Each prior paper sits on one of them. The symbols are those of the note:
$\sigma=\Lambda^2$, $\ell_0=m_Q/\sigma$, $\gst,\bst$ the end velocity in the pair rest frame, $\rho=\gst\bst$,
$\lmax=2(\gst-1)\ell_0$ the oscillation amplitude, and the dimensionless field strengths
\begin{equation}
  \epsB=\frac{qB}{\sigma}=2.254\times10^{-4}\,\frac{B}{3.8112\tesla}\Big(\frac{\keV}{\Lam}\Big)^2,\qquad
  \epsp=\frac{q|\Estar_\perp|}{\sigma},\qquad \Estar=\gamma_{\rm cm}\boldsymbol\beta_{\rm cm}\times\mathbf B . \tag{note 3}
\end{equation}

\paragraph{Step 1: field $\to$ torque $\to$ angular momentum.} In the pair rest frame the lab field $\mathbf B$ appears as
a weak electric field $\Estar$. The magnetic forces on the two ends cancel in the relative motion, so only $\Estar$
exerts a torque, $d\mathbf J/dt=q(\mathbf x_1-\mathbf x_2)\times\Estar$. The separation reverses every half period, so
$\mathbf J$ builds up for one half period and is returned in the next. It is zero at production and at even crossings, and maximal at odd ones.
\emph{Prior work:} Knapen et al.\ Eqs.~(5)--(7); Li et al.\ 2020 Eqs.~(A.4)--(A.5).

\paragraph{Step 2: angular momentum $\to$ miss distance.} At a crossing the string has zero length and $\mathbf J=b\,p^*$, so the ends miss at odd
crossings by
\begin{equation}
  \bodd = 2\,\epsp\,\lzero\Big[\gst-\frac{\operatorname{asinh}\rho}{\rho}\Big]
  \ \xrightarrow{\ \rm NR\ }\ \frac43\,\frac{q|\Estar_\perp|\,m_Q\,\beta^{*2}}{\sigma^2},
  \qquad b_{\rm even}=0, \tag{note 5}
\end{equation}
with a median of about $3.5\um\,(\keV/\Lam)^4$ for $m_Q=500\GeV$ at 3.8\,T.
\emph{Prior work:} Kang--Luty Eq.~(5.64); Li et al.\ 2020 Eq.~(V.8); Li et al.\ 2019 Eq.~(III.12).

\paragraph{Step 3: near miss $\to$ logarithm $\to$ rotation per period (the new step).} Near an odd crossing the
string pulls the ends sideways with $\sigma\bodd/|x|$. Integrated over $\bodd\ll|x|\ll\lmax$ this gives a logarithm. Because the string carries no
momentum at the crossing, the resulting kick rotates all later motion rigidly. The oscillation axis turns toward $+\Estar_\perp$ by
\begin{equation}
  \rot=\epsp\,G(\gst)\Big[\ln\frac{\lmax}{\bodd}+K(\gst)\Big],\qquad
  G=\frac{2[\gst-\operatorname{asinh}\rho/\rho]}{\rho\,\operatorname{artanh}\bst}\xrightarrow{\rm NR}\frac43,\qquad
  K_{\rm NR}=3\ln2-2 , \tag{note 7, 8}
\end{equation}
once per period, always in the same direction, so that $\Theta=N_{\rm odd}\rot$.
\emph{Prior work:} Li et al.\ 2020's $\eta$ (Eq.~A.6), determined numerically; Knapen et al.'s $\Delta\phi$ (Eqs.~8--9),
estimated without the logarithm.

\paragraph{Step 4: accumulation $\to$ plane thickness.} The plane fit (smallest eigenvalue of $T_{ij}$, note~(4)) sees the
out-of-plane part of the accumulated rotation:
\begin{equation}
  \DsB=\frac{qBR}{\sigma}\,F\Big[\ln\frac1{\epsB}+C\Big],\qquad
  c=\frac{1}{\ln(1/\epsB)+C}\simeq0.104,\qquad \DsB\propto B^{1-c}\Lam^{-2+2c}\approx B^{0.9}\Lam^{-1.8}. \tag{note 9, 10}
\end{equation}
The thickness is nearly flat in $m_Q$: $33.5/31.1/29.5\um$ at $200/500/1000\GeV$ and $1\keV$. The blind numerics give
$c=0.105\pm0.002$ and a $B$ exponent of $0.911\pm0.003$ (prediction 0.914).
\emph{Prior work:} the observable is Knapen et al.\ Eqs.~(10)--(11); the secular structure is Li et al.\ 2020 Eq.~(V.5).

\begin{table}[htbp]
\centering\small
\caption{Where each prior paper sits on the P1 chain.}
\label{tab:chain}
\begin{tabularx}{\linewidth}{>{\raggedright\arraybackslash}p{3.2cm}>{\raggedright\arraybackslash}X>{\raggedright\arraybackslash}p{3.4cm}}
\toprule
step & prior statement & P1 equation\\
\midrule
1 torque, $\mathbf J$ & Knapen (5)--(7); Li20 (A.4)--(A.5) & note \S3.2\\
2 miss distance & Kang--Luty (5.64), (6.2); Li20 (V.8); Li19 (III.12) & note (5)\\
3 rotation per period & Li20 $\eta$ (A.6), fitted numerically, (A.17); Knapen $\Delta\phi$ (8)--(9) & note (6)--(8)\\
4 thickness & Knapen $\Delta s$ (10)--(11); Li20 (V.5)--(V.6), (B.9) & note (9)--(10)\\
\bottomrule
\end{tabularx}
\end{table}

%======================================================================================================
\section{Kang \& Luty, arXiv:0805.4642 (JHEP 11 (2009) 065)}\label{sec:kl}

\subsection{What they wrote}
\noindent\textbf{Equation of motion and the straight string} (\S3.2, p.~11--12):
\begin{srcquote}
\[
  \frac{\partial}{\partial t}(m\gamma\vec v)=-\sigma\Big[\sqrt{1-\vec v_\perp^{\,2}}\,\hat s
  +\frac{v_\parallel}{\sqrt{1-\vec v_\perp^{\,2}}}\,\vec v_\perp\Big]+\vec F_{\rm ext}\tag{KL 3.8}
\]
``The gauge string is a dynamical object with its own complicated equation of motion. However, if the quirks have no
further interactions after they are produced (e.g.\ with matter in the detector) then in the center of mass frame the
string remains straight. [\dots] If the string force is much larger than matter forces $F_{\rm ext}\ll\Lambda^2$ (3.10)
then we expect that the string will remain approximately straight in the center of mass frame. [\dots] The full string
dynamics is sufficiently complicated that it would be useful to check this by direct simulation.''
\src{Kang \& Luty, \S3.2, Eqs.~(3.8)--(3.10) and following text.}
\end{srcquote}

\noindent\textbf{The magnetic field} (\S5.8, p.~36):
\begin{srcquote}
``Another effect that can be important in preventing annihilation is the magnetic field in the detector [\dots]. The
quirk center of mass frame is boosted relative to the lab frame, so there will be an electric field in this frame.
This electric field will typically have a component perpendicular to the direction of the quirk motion, which will
give rise to a repulsive force between oppositely charged quirks. This in turn will generate a classical separation
after one oscillation of order
\[
  \Delta r\sim aT^2\sim\frac{v_{\rm cm}B}{m_Q}\Big(\frac{m_Q}{\Lambda^2}\Big)^2 .\tag{KL 5.64}
\]
where $v_{\rm cm}$ is the velocity of the center of mass frame. Demanding that this is larger than $m_Q^{-1}$ (and assuming
$v_{\rm cm}\sim1$) gives $\Lambda\lesssim{\rm MeV}\,(B/{\rm Tesla})^{1/4}(m_Q/{\rm TeV})^{1/2}$ (5.65) [\dots]
It is straightforward to check that the magnetic field in the center of mass frame does not cause a transverse quirk
separation to leading order in the magnetic field.''
\src{Kang \& Luty, \S5.8, Eqs.~(5.64)--(5.65). The same estimate is repeated in \S6, Eq.~(6.2), for mesoscopic strings.}
\end{srcquote}

\subsection{Assessment}
\begin{verdict}
\vhead{Correct in scale, used for a different purpose, and slightly incomplete.}
\begin{enumerate}[leftmargin=*]
\item \textbf{The scale is right.} With $|\Estar_\perp|=\gamma_{\rm cm}\beta_{\rm cm}B\times(\text{angle factor}\le1)$
  and $\sigma=\Lambda^2$, the non-relativistic limit of note~(5) reads
  \[
    \bodd\simeq\frac43\,\beta^{*2}\,\gamma_{\rm cm}\beta_{\rm cm}\,\frac{qB\,m_Q}{\Lambda^4}\times(\text{angle factor}),
  \]
  which is (KL~5.64) with explicit $O(1)$ factors: $\tfrac43\beta^{*2}$, the angle factor, and $\gamma_{\rm cm}$.
  The note's \S3.2 says this explicitly and credits Eq.~(5.64) for the $m_QB/\Lambda^4$ scale.
\item \textbf{Refinement: the separation is not permanent.} ``A separation after one oscillation'' holds at
  every \emph{odd} crossing (one half period after the last even one). At even crossings $\mathbf J$ has returned to zero
  and the ends meet head-on again at linear order in the field (step~1). For Kang--Luty's question (re-annihilation) this
  is not essential, since at an even crossing the axis has already rotated and matter and the next odd crossing intervene. It is
  essential for the thickness, because the logarithm of step~3 arises only at odd crossings.
\item \textbf{The CM magnetic field.} Their statement that $\Bstar$ causes no transverse separation at leading order agrees
  with P1 (note \S3.7: in linear response $\Bstar$ neither changes $\bodd$ nor rotates the axis). P1 finds that $\Bstar$
  does act \emph{non-linearly} inside the odd-crossing region ($-7$ to $-21\%$ on $\rot$ at $\gst=2$, $\kB=4$--16; \EST{}
  in the note).
\item \textbf{The straight-string check they asked for.} Their ``it would be useful to check this by direct simulation''
  is answered in note \S3.7. A step-converged straight-string integrator reproduces the $B\ln(1/B)$ response ($c$ within
  $\approx3\%$ at $\gst=1.28$), but the size of the rotation differs. The predicted straight/exact thickness ratio at $\gst=2$, 0.50, fails:
  the measurement is $0.75\pm0.08$. Evans \& Luty (arXiv:1811.08903, \S2) made the same assumption in the pair rest frame:
  ``As the quirks begin with no relative angular momentum and typically acquire very little relative to the overall
  excitation of the system, retarded time corrections can also be safely neglected.'' P1 finds the retardation factor
  $1/\operatorname{artanh}\bst$ instead of $1/\bst$ in the crossing force (note~(6)); it matters at the level of
  the coefficient $G$, not of the structure.
\end{enumerate}
\textbf{Where P1 enters:} note~(5) is the exact relativistic form of (KL~5.64), restricted to odd crossings. Kang--Luty
make no statement about a plane thickness, a logarithm or a rotation.
\end{verdict}

%======================================================================================================
\section{Knapen, Lou, Papucci \& Setford, arXiv:1708.02243 (PRD 96 (2017) 115015)}\label{sec:knapen}

\subsection{What they wrote}
\noindent\textbf{The torque estimate} (\S II, p.~3):
\begin{srcquote}
``To estimate the effect of the B-field, it is useful to move to the center of mass frame. In this frame, the magnetic
field is seen as an E-field, which introduces a torque
\[
  \vec\tau\sim2\vec d\times(e\vec E_{\rm cm})=2e\gamma_{\rm cm}\,\vec d\times(\vec v_{\rm cm}\times\vec B),\tag{Kn 5}
\]
[\dots] 2d is the typical displacement of the quirks in the center of mass frame, and B the magnetic field in the lab
frame. The angular momentum that is picked up in a single oscillation with period $\Delta t$ is roughly
\[
  L\sim|\tau|\Delta t\sim e\,v_{\rm cm}\gamma_{\rm cm}\,v\gamma(\gamma-1)\,\frac{m_Q^2B}{\Lambda^4}\tag{Kn 6}
\]
[\dots] While the internal angular momentum of the system is typically very large in units of $\hbar$, it is still small
compared to the resolution of the trackers, and the trajectories remain co-planar as far as the experiments are
concerned. We can see this by estimating the angular rotation of the plane spanned by the quirks' velocity vectors:
\[
  \Delta\phi\sim|\tau|\Delta t^2/I\sim e\,v_{\rm cm}\gamma_{\rm cm}\,\frac{\gamma^2v^2B}{(\gamma-1)^2\Lambda^2}\tag{Kn 8}
\]
\[
  \sim10^{-5}\,\frac{v_{\rm cm}}{v^2}\Big(\frac{B}{2\,{\rm T}}\Big)\Big(\frac{2\,{\rm keV}}{\Lambda}\Big)^2 .\tag{Kn 9}
\]
The key point here is that the effect of the torque on the angular acceleration is suppressed by the large moment of
inertia of the system $I\sim2d^2m_Q$. There could be an enhancement for close to threshold quirks, where $v\ll1$; but
this is relevant only for a very small part of phase space, and $\Delta\phi$ is typically not larger than $10^{-3}$. We
show $\Delta\phi$ in Fig.~3 for two example events, as found in the full numerical solution of Eq.~3 with $F_{\rm ext}$
the Lorentz force. In the numerical result, $\Delta\phi$ oscillates and slowly accumulates as the quirks travel through
the detector until it stabilizes around a fixed value. We see that the typical $\Delta\phi$ is somewhat larger than
$10^{-5}$, but is still small compared to the resolution of the tracker. The effect of the magnetic field is accounted
for in all our simulations, and any potential efficiency loss due to shifting of the quirks' plane is included in our
results.''
\src{Knapen et al., \S II, Eqs.~(5)--(9) and following text. Fig.~3 (y-axis in units of $10^{-3}$) shows the accumulated
$\Delta\phi\approx0.4\times10^{-3}$ for $m_Q=800\GeV$ and $\approx0.1\times10^{-3}$ for $1.8\TeV$, at $\Lambda=2\keV$, $B=2$\,T.}
\end{srcquote}

\noindent\textbf{Material and the scope of the argument} (\S II):
\begin{srcquote}
``[\dots] the forces exerted by the ionization process when the quirks traverse a silicon layer are of the order
$\sim(100\ {\rm eV})^2$ [\dots] such that the shift in angle is $\Delta\phi\sim10^{-3}$ for each layer the quirks traverse.
We therefore neglect this effect in our simulations.''
\end{srcquote}

\noindent\textbf{The observable} (\S III.C):
\begin{srcquote}
\[
  d(\mathbf n,\mathbf x^a)\equiv\sqrt{\frac{1}{N-1}\sum_{a=1}^N(\mathbf n\cdot\mathbf x^a)^2}\,,\qquad
  T(\mathbf x^a)_{ij}\equiv\frac{1}{N-1}\sum_{a=1}^Nx^a_ix^a_j \tag{Kn 10, 11}
\]
``Minimization of $d$ with respect to $\mathbf n$ simply reduces to solving an eigenvalue problem for $T$. The smallest
eigenvalue, $\Delta s^2$, then gives the minimum value of $d^2$ [\dots]. $\Delta s$ therefore gives a measure of the
thickness of the plane.'' Their selection uses $\Delta s<0.05$\,cm in the plane finding and $\Delta s<0.01$\,cm as a
signal cut.
\end{srcquote}

\subsection{Assessment}
\begin{verdict}
\vhead{The scale and the $m_Q$ independence are right; the secular, logarithmic term is missing.}
\begin{enumerate}[leftmargin=*]
\item \textbf{What (Kn~8) captures.} With $\tau\sim2d\,qE$, $I\sim2d^2m_Q$, $d\sim(\gamma-1)m_Q/\Lambda^2$ and
  $\Delta t\sim2v\gamma\,m_Q/\Lambda^2$, the estimate is $\Delta\phi\sim qE\,\Delta t^2/(d\,m_Q)\propto qE/\sigma$, i.e.\
  $\epsp$ times a velocity factor. The masses cancel. This is the right scale, and per period the P1 rotation (note~(7)) is
  also $\propto\epsp$ and independent of $m_Q$ at fixed velocities. The note credits Knapen et al.\ for the scale and
  the $m_Q$-independent $B/\Lambda^2$ behaviour (\S3.4 ``Credit'', \S3.5).
\item \textbf{What it misses.} The estimate uses torque $\times\,\Delta t^2/I$ \emph{within} one oscillation. That is the
  back-and-forth swing that $\mathbf J$ makes inside each period (step~1), which is undone in the next half period. It
  does not identify the part that survives: the kick at each odd crossing, which carries $\ln(\lmax/\bodd)\approx8$--14
  and always points toward $+\Estar_\perp$, so that it adds coherently over $N_{\rm odd}$ periods (note~(7)). The note's
  wording is that the estimate ``captures the scale $q|\Estar|/\sigma$ and the alternating half-period rotation, but not
  the logarithmic secular term.''
\item \textbf{Consequence for the scaling.} Because of the logarithm the response is $B\ln(1/B)$, not $B$. The
  thickness exponents are $B^{0.9}\Lambda^{-1.8}$, not $B\,\Lambda^{-2}$ (note~(10); measured $B$ exponent
  $0.911\pm0.003$).
\item \textbf{Their numerics vs their estimate} (\EST{}, our reading). They note that the numerical $\Delta\phi$ is
  ``somewhat larger than $10^{-5}$'' and their Fig.~3 shows $(1$--$4)\times10^{-4}$. Two remarks:
  (a) In P1 the accumulated angle after a distance $R$ is $\Theta=N_{\rm odd}\rot$ with $N_{\rm odd}\propto\sigma R/m_Q$,
  so $\Theta\propto qBR/p^*$, a bending angle (note \S3.5). The accumulated angle therefore depends on $m_Q$ even though the
  rotation per period does not, and the thickness (angle $\times$ lever arm $\propto m_Q$) is flat in $m_Q$. This is
  consistent with Fig.~3, where the 800\,GeV event has a larger $\Delta\phi$ than the 1.8\,TeV one.
  (b) We do not draw a size comparison between note~(7) and (Kn~9). The project's novelty audit (erratum, 11 Oct) found
  that evaluating (Kn~8) gives a value about $9\times$ larger than the printed $10^{-5}$ of (Kn~9), so the ``estimate'' is
  itself uncertain by an order of magnitude. Their statement that $\Delta\phi$ ``stabilizes around a fixed value'' was
  not analysed in the project.
\item \textbf{The observable.} P1 predicts exactly their $\Delta s$ (note~(4) is (Kn~10--11)). At $1\keV$ the field part
  is $\DsB\approx30\um$, well inside their $\Delta s<0.01$\,cm cut. It grows to $\sim80\um$ near $0.6\keV$ (\EST{},
  from $\Lambda^{-1.8}$), which is where a fixed 100\,$\mu$m cut starts to cost efficiency. They simulated the field
  numerically (``any potential efficiency loss \dots{} is included in our results''), so \textbf{their reach is
  unaffected}.
\end{enumerate}
\textbf{Where P1 enters:} note~(7) replaces (Kn~8--9) by an analytic rotation per period with the logarithm, and
note~(9) predicts their $\Delta s$ in closed form.
\end{verdict}

%======================================================================================================
\section{Li, Li, Pei \& Zhang 2019, arXiv:1911.02223 (PRD 102 (2020) 056006)}\label{sec:li19}

\subsection{What they wrote}
\begin{srcquote}
``Given a set of hit positions $\vec h_i$ ($i=1,2,\dots,N$), we can define
\[
  d(\vec n)=\sqrt{\frac{1}{N-1}\sum_{i=1}^N(\vec n\cdot\vec h_i)^2}\tag{Li19 III.11}
\]
[\dots] The $\vec n$ that gives the smallest $d(\vec n)$ (denoted by $d_{\min}$) corresponds to the normal vector of the
quirk plane. The thickness of the quirk plane is induced by the $\vec n$-component of the Lorentz force, so it can be
estimated as [50]
\[
  d_{\min}\sim1.16\,\Big(\frac{m}{100\ {\rm GeV}}\Big)\Big(\frac{\rm keV}{\Lambda}\Big)^4\Big(\frac{|q|}{e}\Big)
  \Big(\frac{B_{xy}}{\rm T}\Big)\ \mu{\rm m},\tag{Li19 III.12}
\]
where $B_{xy}=|\vec B-(\vec B\cdot\vec n)\vec n|$. The thickness of the quirk plane can also be calculated by the
algorithm proposed in Ref.~[1], which is simplified to an eigenvalue problem. We have verified that the thickness
calculated from this method matches well with our estimation in Eq.~(III.12). In the lower panel in Fig.~4, we plot the
distributions of the quirk plane thickness, varying both the quirk mass and the confinement scale. [\dots] It is clear
that the dependence on the confinement scale is much stronger than the dependence on quirk mass. We can also observe
that for $\Lambda\in[O(100)\ {\rm eV},O(1)\ {\rm keV}]$, most of the events have a quirk plane thickness smaller than
$\sim O(100)\ \mu$m.''
\src{Li et al.\ 2019, \S III.C, p.~5. Their Ref.~[50] is Li et al.\ 2020 (2002.07503); Ref.~[1] is Knapen et al.}
\end{srcquote}

\subsection{Assessment}
\begin{verdict}
\vhead{Eq.~(III.12) is the miss distance, not the plane thickness.}
\begin{enumerate}[leftmargin=*]
\item \textbf{The scaling identifies it.} (Li19~III.12) $\propto m_Q\,qB/\Lambda^4$ is the scaling of step~2, the miss
  distance $\bodd$ (note~(5); Kang--Luty~5.64). The accumulated plane-fit thickness of step~4 scales as
  $\approx B^{0.9}\Lambda^{-1.8}$ and is nearly flat in $m_Q$. The note's \S3.2 states this directly: ``it measures the miss distance, not the accumulated plane-fit thickness, which scales differently.''
\item \textbf{A tension inside their paper.} (Li19~III.12) is linear in $m$, while their Fig.~4 (lower panel), computed
  numerically, shows that ``the dependence on the confinement scale is much stronger than the dependence on quirk mass''.
  The numerics agree with P1, not with (III.12).
\item \textbf{Numbers.} Table~\ref{tab:li19} compares (III.12) with the P1 median thickness and the measured median at
  $m_Q=500\GeV$. With $B_{xy}=B$ (an upper bound on (III.12)) the formula is $\approx6.3\times$ the median $\bodd$ at every
  $\Lambda$, i.e.\ it has exactly the $\bodd$ scaling. It agrees with the true thickness in order of magnitude near
  $1\keV$, by coincidence of scales (\EST{}, our reading of why they found it to ``match well''), and falls below it
  by a factor that grows as $\Lambda^{2.2}$.
\item \textbf{Mass.} From 200 to 1000\,GeV (III.12) grows $\times5$; the P1 model medians at 1\,keV are
  $33.5/31.1/29.5\um$, i.e.\ they fall by 12\%.
\item \textbf{Practical caveat.} Above $\Lcross\approx1.2$--$1.9\keV$ multiple scattering, not the field, sets the
  observed thickness (floor $10$--$28\um$; note \S4). The large factors at 5--10\,keV in Table~\ref{tab:li19} concern
  the field part only.
\end{enumerate}
\textbf{Where P1 enters:} it reinterprets (III.12) as the odd-crossing miss distance and replaces it, as a thickness
estimate, by note~(9). Their collider analysis uses numerically simulated hits, so \textbf{their reach is unaffected}; the
formula is misleading only if it is used to estimate a thickness.
\end{verdict}

\begin{table}[htbp]
\centering\small
\caption{Eq.~(III.12) of Li et al.\ 2019 against P1, $m_Q=500\GeV$, $B=3.8112$\,T, $|q|=1$. (III.12) is evaluated with
$B_{xy}=B$, its largest value. Median $\bodd\approx3.5\um\,(\keV/\Lam)^4$ (note \S3.2). P1 model and blind-numerics medians
of $\DsB$ from the Stage-1 note (\rep{analytic/NEAR\_MISS\_IMPULSE.md}, verification table row 20), field only, no material.}
\label{tab:li19}
\begin{tabular}{rrrrrr}
\toprule
$\Lambda$ [keV] & (III.12) [$\mu$m] & median $\bodd$ [$\mu$m] & P1 model $\DsB$ [$\mu$m] & blind numerics [$\mu$m] & model / (III.12)\\
\midrule
1  & 22.1    & 3.5      & 31.1 & $29.1\pm2.1$  & 1.4\\
2  & 1.38    & 0.22     & 8.6  & $8.17\pm0.72$ & 6.2\\
3  & 0.273   & 0.043    & 4.18 & $3.80\pm0.29$ & 15\\
5  & 0.0354  & 0.0056   & 1.61 & $1.55\pm0.11$ & 45\\
10 & 0.00221 & 0.00035  & 0.45 & $0.46\pm0.06$ & 200\\
\bottomrule
\end{tabular}
\end{table}

\paragraph{A related finding not in the note.} During the theory-synthesis review (8 Oct, \rep{theory\_track/LOG.md}),
the critical reviewer coded the literal Li et al.\ integrator (Li19 Eqs.~III.6--III.10). Step-converged, it reproduces
the exact string's thickness to 0.84--1.04 (3--35\% low depending on $\gst$). At the step size printed in Eq.~(III.10)
it is not converged, and the thickness changes by $\times0.8$--20. This is a project-internal result that was not carried into
the note and has not been re-checked here (\EST{}).

%======================================================================================================
\section{Li, Li, Pei \& Zhang 2020, ``The quirk trajectory'', arXiv:2002.07503 (EPJC 80 (2020) 651)}\label{sec:li20}

This is the closest prior work. P1 can be described as an analytic theory of their rotation parameter $\eta$.

\emph{Notation of Li et al.:} in the CoM frame the quirks move along $\pm\hat e_{xc}$; $z_{1,2}$ are their displacements
along $\hat e_z$, the direction of the field-induced force; $z^\pm=\tfrac12(z_1\pm z_2)$. The variable $g$ counts half
periods: $g=0.5$ is the first turning point, $g=1$ the first (odd) crossing, $g=2n$ the even crossings. $r(g)$ is the
distance of a quirk from the CoM, $r(0.5)=\lmax/2$; $\rho=\gst\bst$ as in the note.

\subsection{What they wrote}
\noindent\textbf{The rotation per period} (App.~A):
\begin{srcquote}
``[\dots] from Eq.~II.12, we know the average torque on the system of two quirks is zero when $g$ is increased from $2n$
to $2n+2$ such that the total angular momentum of the quirk pair is also invariant at $g=2n$. Then there must be
\[
  z^+(2n)=0,\qquad z_1(2n)=-z_2(2n)=z^-(2n)=0,\qquad n=0,1,2\dots\tag{Li20 A.5}
\]
This means that two quirks meet each other and have opposite velocity at $g=2n$ in the CoM frame. We hence conclude that
kinematics of quirk system at $g+2$ can be obtained by rotating the system at $g$ with an angle $-\eta$ around $\hat e_{yc}$,
leading to
\[
  z^-(2n+\Delta)=z^-(\Delta)\cos(n\eta)+(-1)^{[\Delta]}r(\bar\Delta)\sin(n\eta)\approx z^-(\Delta)+(-1)^{[\Delta]}r(\Delta)\,n\eta\,,\tag{Li20 A.6}
\]
where $\Delta\in[0,2]$ and $\sin\eta\approx\eta\ll1$.''
\end{srcquote}

\noindent\textbf{First-order theory fails near the crossing} (App.~A, after Eq.~A.13):
\begin{srcquote}
``Results from the numerical calculations show that the difference between $z^-_o(g)$ and $z^-_{o0}(g)$ can not be
ignored when $g\in[0.5,1]$, even though they are still of the same order of magnitude. Using $z^-(0.5)=w_{0.5}$,
$z^-(1)=w_1$ from the numerical calculations, we can approximately construct $z^-_o(g)$ [\dots]''
\end{srcquote}

\noindent\textbf{$\eta$ is not derived} (App.~A):
\begin{srcquote}
``The above unknown parameters $w_{0.5}$, $w_1$ and $\eta$ can be obtained from the numerical calculations. Moreover, the
magnitude order relation
\[
  r(0.5)\,|\eta|\sim|w_1|\sim|z^-_{o0}(1)|,\tag{Li20 A.17}
\]
is also obtained by the numerical results.''
\end{srcquote}

\noindent\textbf{The thickness formula} (\S V):
\begin{srcquote}
``According to the discussions in Sec.~II~B and Appendix~A, the quirk pair plane in the CoM frame can be approximately
obtained by rotating the $\hat e_{xc}$--$\hat e_{yc}$ plane by an angle $(-j\eta)$ around $\hat e_{yc}$ [\dots] thickness
of quirk pair plane in the lab frame can be estimated as
\[
  d=\sqrt{\frac{\eta^2r^2(0.5)}{3l_0}\sum_{i=1}^{l_0}\Big(j(l_0)-\Big[\frac i2\Big]\Big)^2
  +\frac43\Big(\frac{2z^{+2}_o(1)}{15}+\frac{z^{-2}_o(1)}{5}\Big)}\tag{Li20 V.5}
\]
\[
  \approx\sqrt{\Big(\frac{1}{3l_0}\sum_{i=1}^{l_0}\Big(j(l_0)-\Big[\frac i2\Big]\Big)^2+\frac4{15}\Big)d_E^2+\frac8{45}d_B^2}\,,\tag{Li20 V.6}
\]
where we have used the relations of Eq.~A.17 in the second line, and [\dots]
\[
  d_E=z^-_{o0}(1)=\frac{2mqE_z}{\Lambda^4}\,\rho\sinh^{-1}[\rho],\qquad
  d_B=z^+_o(1)=\frac{2mqB_c}{\Lambda^4}\big(\rho-\tan^{-1}[\rho]\big).\tag{Li20 V.8, V.9}
\]
We provide the validation of Eq.~V.6 in Appendix~C, where we can conclude that our analytic formula for the quirk pair
plane thickness matches the numerical result within an order magnitude when quirk pair goes through $\sim O(1-100)$
periods inside tracker. The main difference [\dots] attributes to the following reasons. First of all, the former uses
the shapes of overall quirk trajectories [\dots]. The latter only uses $O(10)$ positions of hits [\dots]. Besides, we
employ the relations of Eq.~A.17 to obtain Eq.~V.6 as an approximation of Eq.~V.5.''
\end{srcquote}

\noindent\textbf{Validation} (App.~C):
\begin{srcquote}
``Strictly speaking, $E_z$ and $B_c$ respectively in Eqs.~V.8 and V.9 change with the number of periods due to the
rotation of the quirk-pair system described by Eq.~A.10, but are thought to be invariant because of the small $\eta$. The
effect of the system rotation can not be ignored when $n_T>O(100)$, so the plane thickness given by Eq.~B.9 is less than
the numerical result.''
\end{srcquote}

\noindent\textbf{Charge dependence} (\S V.A) and \textbf{fixed string direction} (\S VIII):
\begin{srcquote}
``Because it is the Lorentz force rendering the quirk traveling outside the plane, the plane thickness is linearly
proportional to each of the quirk charges. This feature is clearly shown by the solid lines in the right panel of
Fig.~4 [\dots]''

``[\dots] we plot the changes in the quirk pair plane thickness when the fixed $\hat s_{1,2}$ are used instead of
Eqs.~VIII.3 and VIII.4. The quirk pair plane thickness is changed dramatically [\dots]. So choosing the correct $\hat s$
direction is critical in coplanar quirk search.''
\end{srcquote}

\subsection{Assessment}
\begin{verdict}
\vhead{The structure is right; $\eta$ is derived in P1; (A.17) drops a logarithm; ``linear in $q$'' holds to $\approx7\%$.}
\begin{enumerate}[leftmargin=*]
\item \textbf{Agreement on the structure.} Their (A.5) is step~1 of P1: $\mathbf J$ returns to zero and the ends meet
  head-on at even crossings. Their (A.6) and the first term of (V.5) are steps~3--4: a small rotation per period that
  accumulates as $n\eta$ and dominates the thickness. The note credits this as the closest prior work (\S3.4 ``Credit'').
\item \textbf{$\eta$ is derived in P1.} Their $\eta$ is the note's $\rot$:
  \[
    \eta\equiv\rot=\epsp\,G(\gst)\Big[\ln\frac{\lmax}{\bodd}+K(\gst)\Big]\qquad\text{(note 7)},
  \]
  ``Equation~(7) is an analytic expression for that angle'' (note \S3.4). It is generated at the odd crossings, which Li et al.\
  do not identify as the origin.
\item \textbf{(A.17) is not a constant ratio} (\EST{}, our arithmetic from note (5) and (7)). With $r(0.5)=\lmax/2$ and
  the exact miss distance $w_1=\bodd/2$,
  \[
    \frac{r(0.5)\,\eta}{\bodd/2}=\frac{2(\gst-1)}{\rho\,\operatorname{artanh}\bst}\Big[\ln\frac{\lmax}{\bodd}+K\Big]
    \simeq(0.88\text{--}1.0)\times\Big[\ln\frac{\lmax}{\bodd}+K\Big]\approx8\text{--}14 .
  \]
  The relation $r(0.5)|\eta|\sim|w_1|$ therefore hides the logarithm, which depends on $B$ and $\Lambda$. Replacing $\eta$
  by $d_E/r(0.5)$ in (V.5)~$\to$~(V.6) turns the true $B\ln(1/B)$ into a linear $B$. Their $d_E$ is itself
  $2\rho\sinh^{-1}\rho/(\gst-\sinh^{-1}\rho/\rho)=3.0$--$4.2$ times $\bodd/2$, because the first-order $z^-_{o0}$ integrates
  the field force without the string's transverse pull (Stage-1 note \S3.2). The two errors partly compensate: $r(0.5)\eta/d_E\approx2$--5.
  This plausibly accounts for part of their ``within an order of magnitude'' (App.~C), but the note does not quantify it.
\item \textbf{Why first order fails on $g\in[0.5,1]$.} They observed this numerically and patched it with the fitted $w_{0.5}$,
  $w_1$. P1 explains it (note \S3.3): near a crossing $|x|\simeq2v|s|$, so the string pull $\sigma z/|x|\propto1/|s|$ and its
  impulse diverges logarithmically unless $z=0$ there. $z=0$ at production and at even crossings, so first-order theory is
  exact there and fails only at odd crossings. The exact Nambu--Goto solution in the intermediate region (note~(6)) resolves the
  divergence.
\item \textbf{Charge.} ``Linearly proportional to each of the quirk charges'' is approximate. Since $q$ enters only via
  $qB$, $\DsB(q)/\DsB(1)=q\,[1-\ln q/(\ln(1/\epsB)+C)]$, i.e.\ $1.86$ instead of 2 for $|q|=2$. This is verified by the blind numerics
  (ratio to prediction $0.999\pm0.001$, note Table~2). The note calls the $\approx c\ln2\approx7\%$ deviation ``consistent
  with their `linearly proportional' at their precision'' (\S3.7).
\item \textbf{Fixed $\hat s$} (\EST{}, from the project's novelty audit, not in the note). A fixed string direction removes the
  near-miss swing of $\hat s$ at the odd crossings, which is the P1 mechanism. Their Fig.~6 (right) is therefore indirect
  numerical evidence that this swing controls the thickness. Li et al.\ do not interpret it this way.
\item \textbf{Their numerics stand.} The published reach uses the numerical thickness, and a step-converged
  straight-string (``Li-method'') integrator does contain the logarithm ($c$ within $\approx3\%$ of the exact string at
  $\gst=1.28$; note \S3.7). The note says this ``without implying an error in published limits''.
\end{enumerate}
\textbf{Where P1 enters:} insert note~(7) for $\eta$ in the first term of (Li20~V.5), instead of the numerical fit and
(A.17). With the coherent accumulation $\Theta=N_{\rm odd}\rot$ and the rotating-yo-yo plane fit this gives note~(9), which
reproduces the exact string to 1--4\% per event for $\gst\lesssim1.6$.
\end{verdict}

%======================================================================================================
\section{Summary}\label{sec:summary}

\begin{table}[H]
\centering\small
\caption{Verdicts. ``Reach affected'' refers to the paper's published collider sensitivity.}
\label{tab:summary}
\begin{tabularx}{\linewidth}{>{\raggedright\arraybackslash}p{2.3cm}>{\raggedright\arraybackslash}X>{\raggedright\arraybackslash}X>{\raggedright\arraybackslash}p{2.4cm}c}
\toprule
paper & their statement & P1 verdict & P1 equation & reach affected\\
\midrule
Kang--Luty (5.64) & separation $\sim(v_{\rm cm}B/m)(m/\Lambda^2)^2$ ``after one oscillation'' &
  right scale; it is $\bodd$, present only at odd crossings & note (5) & no (none given)\\
Knapen (8)--(9) & $\Delta\phi\propto B/\Lambda^2$, $m_Q$-independent & scale and $m_Q$ independence right;
  coherent log term missing; $B\ln(1/B)$, $\DsB\approx B^{0.9}\Lam^{-1.8}$ & note (7), (9) & no\\
Li 2019 (III.12) & thickness $\propto mB/\Lambda^4$ & \textbf{misidentified}: miss distance, not thickness;
  $\times6$--200 low at 2--10\,keV & note (5) vs (9) & no\\
Li 2020 (A.6), (A.17), (V.5) & rotation $\eta$ per period, numerical; $r\eta\sim d_E$ & structure right; $\eta$
  derived; (A.17) drops $\ln(\lmax/\bodd)\approx8$--14 & note (7) & no\\
Li 2020 \S V.A & thickness linear in $q$ & approximate: 1.86 for $|q|=2$ & note \S3.7 & no\\
Kang--Luty \S3.2; Evans--Luty \S2 & straight string, retardation negligible & log structure survives; coefficient
  differs ($1/\operatorname{artanh}\bst$); $\gst=2$ ratio prediction fails & note (6), \S3.7 & no\\
\bottomrule
\end{tabularx}
\end{table}

\paragraph{Where to read this in the note.} \S3.2 end (Kang--Luty, Li~V.8, Li~III.12); \S3.3 (first-order breakdown, Li App.~A);
\S3.4 ``Credit'' (Li~$\eta$, Knapen torque); \S3.5 after Table~1 (Knapen and Li scaling); \S3.7 ``Charge and spin'' and
``Straight strings and a failed prediction''; App.~D around the ``alternating pattern'' (Li~A.4--A.5).

\paragraph{References.}
{\small
\begin{enumerate}[label={[\arabic*]},leftmargin=*]
\item J.~Kang, M.~A.~Luty, \emph{Macroscopic strings and `quirks' at colliders}, JHEP 11 (2009) 065, arXiv:0805.4642.
\item S.~Knapen, H.~K.~Lou, M.~Papucci, J.~Setford, \emph{Tracking down quirks at the Large Hadron Collider}, Phys.\ Rev.\ D 96 (2017) 115015, arXiv:1708.02243.
\item J.~Li, T.~Li, J.~Pei, W.~Zhang, \emph{Uncovering quirk signal via energy loss inside tracker}, Phys.\ Rev.\ D 102 (2020) 056006, arXiv:1911.02223.
\item J.~Li, T.~Li, J.~Pei, W.~Zhang, \emph{The quirk trajectory}, Eur.\ Phys.\ J.\ C 80 (2020) 651, arXiv:2002.07503.
\item J.~A.~Evans, M.~A.~Luty, \emph{Stopping quirks at the LHC}, JHEP 06 (2019) 090, arXiv:1811.08903.
\item Our note: \emph{Coplanar, non-helical tracker patterns of electroweak quirks at the HL-LHC}, v0.4, 17 Oct 2026 (\rep{theory\_track/paper/quirks\_p1p2\_note.pdf}); novelty audit \rep{theory\_track/analytic/NOVELTY\_AUDIT.md}; Stage-1 note \rep{theory\_track/analytic/NEAR\_MISS\_IMPULSE.md}.
\end{enumerate}}

\end{document}
